% ============================================================
%  LAMPIRAN 5 - KUERI CYPHER
%  Compile standalone:  pdflatex lampiran-5.tex
%  Compile as part of thesis: pdflatex main.tex
% ============================================================
\documentclass[../../thesis]{subfiles}
\usepackage{hyphenat}
\usepackage{iftex}

% ── Encoding, Language & Fonts ───────────────────────────────
\ifXeTeX
    \usepackage{polyglossia}
    \setmainlanguage{bahasai}
    \usepackage{fontspec}
    \setmainfont{TeX Gyre Termes}
    \setsansfont{TeX Gyre Heros}
    \setmonofont[Scale=0.88]{TeX Gyre Cursor}
\else
    \usepackage[utf8]{inputenc}
    \usepackage[T1]{fontenc}
    \usepackage[indonesian]{babel}
    \usepackage{mathptmx}
    \usepackage{helvet}
    \usepackage{courier}
\fi

% ── Page geometry (A4, UI thesis margins) ────────────────────
\usepackage[
  a4paper,
  top    = 2.54cm,
  bottom = 2.54cm,
  left   = 3.17cm,
  right  = 2.54cm
]{geometry}

% ── Line spacing & paragraph ─────────────────────────────────
\usepackage{setspace}
\onehalfspacing
\setlength{\parindent}{1.27cm}
\setlength{\parskip}{0pt}

% ── Microtype & justification ────────────────────────────────
\usepackage{microtype}
\sloppy

% ── Headings ─────────────────────────────────────────────────
\usepackage{titlesec}

\titleformat{\chapter}[block]
  {\centering\rmfamily\bfseries\large}{}
  {0pt}{\MakeUppercase}
\titlespacing*{\chapter}{0pt}{0pt}{1.5em}

\titleformat{\section}
  {\rmfamily\bfseries\normalsize}{\thesection}{1em}{}
\titlespacing*{\section}{0pt}{1.2em}{0.6em}

\titleformat{\subsection}
  {\rmfamily\bfseries\normalsize}{\thesubsection}{1em}{}
\titlespacing*{\subsection}{0pt}{1em}{0.5em}

% ── Code listings (only used when compiled standalone) ───────
\usepackage{listings}
\lstdefinestyle{promptbox}{
  basicstyle      = \ttfamily\scriptsize,
  breaklines      = true,
  breakatwhitespace = false,
  columns         = fullflexible,
  keepspaces      = true,
  showstringspaces = false,
  frame           = single,
  framesep        = 5pt,
  xleftmargin     = 10pt,
  xrightmargin    = 2pt,
  aboveskip       = 8pt,
  belowskip       = 8pt,
}
\lstdefinestyle{cypherbox}{
  style        = promptbox,
  morecomment  = [l]{//},
  morecomment  = [s]{/*}{*/},
  keywordstyle = \bfseries,
  morekeywords = {MATCH,OPTIONAL,MERGE,CREATE,CONSTRAINT,FOR,REQUIRE,IS,UNIQUE,
                  WHERE,WITH,RETURN,SET,UNWIND,CALL,YIELD,DETACH,DELETE,AS,AND,
                  OR,NOT,IN,CASE,WHEN,THEN,ELSE,END,STARTS,CONTAINS,ORDER,BY,
                  LIMIT,DISTINCT,coalesce,count,sum,avg,collect,toFloat,length,
                  nodes,type,trim,labels,relationshipTypes},
}

% ── Hyperlinks ───────────────────────────────────────────────
\usepackage[hidelinks]{hyperref}

% ============================================================
\begin{document}

% When standalone: switch to appendix mode and set the counter
% so appendix-local numbering renders correctly. When compiled inside main.tex,
% main.tex issues \appendix and advances the counter before \subfile.
\ifSubfilesClassLoaded{}{\appendix\setcounter{chapter}{5}}

\lampiran{Kueri Cypher}
\label{lamp:cypher}

Lampiran ini memuat kueri Cypher yang dipakai untuk (1)~mengonstruksi
\textit{knowledge graph} ke Neo4j dengan strategi \texttt{MERGE} yang idempoten
(Bab~4, Subbab~\ref{sec:graph-completion}) dan (2)~menghitung metrik evaluasi
struktural, ekstraksi, dan hubungan antar-entitas yang definisinya diberikan
pada Bab~2. Kueri konstruksi diparametrikan; \textit{placeholder} bertanda
\texttt{\$nama} disubstitusi dari data hasil ekstraksi--konsensus, dan
\texttt{:\{REL\_TYPE\}} menandai tipe relasi dinamis yang ditentukan saat
eksekusi. Kredensial koneksi Neo4j tidak disertakan.

% ==================================================================
\subsection{Konstruksi Graf (Strategi \texttt{MERGE})}
\label{lamp:cypher-konstruksi}

\subsubsection{Batasan Keunikan (\textit{Constraints})}

Sebelum \textit{ingest}, lima \textit{constraint} keunikan dibuat agar
\texttt{MERGE} bersifat idempoten dan tidak menduplikasi \textit{node}.

\begin{lstlisting}[style=cypherbox,caption={Constraint keunikan node.},label={lst:cypher-constraint}]
CREATE CONSTRAINT consensus_grade_name IF NOT EXISTS
  FOR (g:Grade) REQUIRE g.name IS UNIQUE;
CREATE CONSTRAINT consensus_chapter_key IF NOT EXISTS
  FOR (c:Chapter) REQUIRE (c.grade, c.name) IS UNIQUE;
CREATE CONSTRAINT consensus_subtopic_key IF NOT EXISTS
  FOR (s:Subtopic) REQUIRE (s.grade, s.chapter, s.name) IS UNIQUE;
CREATE CONSTRAINT consensus_concept_key IF NOT EXISTS
  FOR (c:Concept) REQUIRE (c.grade, c.name) IS UNIQUE;
CREATE CONSTRAINT consensus_concept_target_key IF NOT EXISTS
  FOR (t:ConceptTarget) REQUIRE (t.grade, t.name) IS UNIQUE;
\end{lstlisting}

\subsubsection{Pass~1 --- Node Struktural}

Pass~1 membuat hierarki \texttt{Grade}~$\rightarrow$~\texttt{Chapter}~$\rightarrow$~\texttt{Subtopic}~$\rightarrow$~\texttt{Concept}
beserta relasi struktural \texttt{HAS\_CHAPTER}, \texttt{NEXT\_CHAPTER},
\texttt{HAS\_SUBTOPIC}, dan \texttt{HAS\_CONCEPT}.

\begin{lstlisting}[style=cypherbox,caption={Pembuatan node dan relasi struktural (Pass~1).},label={lst:cypher-pass1}]
// Grade
MERGE (g:Grade {name: $name})
SET g.type = $type, g.source = 'consensus', g.graph_role = 'structure';

// Chapter + (Grade)-[:HAS_CHAPTER]->(Chapter)
MERGE (c:Chapter {name: $name, grade: $grade})
SET c.summary = $summary, c.previous = $previous, c.next = $next,
    c.subchapters_json = $subchapters_json,
    c.source = 'consensus', c.graph_role = 'structure'
WITH c
MATCH (g:Grade {name: $grade})
MERGE (g)-[r:HAS_CHAPTER]->(c)
SET r.source = 'consensus', r.graph_role = 'structure';

// Urutan bab: (Chapter)-[:NEXT_CHAPTER]->(Chapter)
MATCH (a:Chapter {name: $source, grade: $grade})
MATCH (b:Chapter {name: $target, grade: $grade})
MERGE (a)-[r:NEXT_CHAPTER]->(b)
SET r.source = 'consensus', r.graph_role = 'structure';

// Subtopic + (Chapter)-[:HAS_SUBTOPIC]->(Subtopic)
MERGE (s:Subtopic {name: $name, chapter: $chapter, grade: $grade})
SET s.source = 'consensus', s.graph_role = 'structure'
WITH s
MATCH (c:Chapter {name: $chapter, grade: $grade})
MERGE (c)-[r:HAS_SUBTOPIC]->(s)
SET r.source = 'consensus', r.graph_role = 'structure';

// Concept + (Subtopic)-[:HAS_CONCEPT]->(Concept)
MERGE (k:Concept {name: $name, grade: $grade})
SET k.description = $description, k.glossary_validated = $glossary_validated,
    k.materi_pokok_ref = $materi_pokok_ref, k.provenance = $provenance,
    k.source = 'consensus', k.graph_role = 'concept'
WITH k
MATCH (s:Subtopic {name: $subtopic, chapter: $chapter, grade: $grade})
MERGE (s)-[r:HAS_CONCEPT]->(k)
SET r.source = 'consensus', r.graph_role = 'structure';
\end{lstlisting}

\subsubsection{Pass~2 --- Relasi Intra-Buku}

Pass~2 menulis relasi semantik antar-konsep dalam satu buku. Jika target
relasi belum berupa \texttt{Concept} yang dikenal, sistem membuat
\textit{node} sementara \texttt{ConceptTarget}. Tipe relasi
(\texttt{:\{REL\_TYPE\}}) bersifat dinamis sesuai \textit{controlled vocabulary}.

\begin{lstlisting}[style=cypherbox,caption={Relasi intra-buku konsep-ke-konsep dan konsep-ke-target (Pass~2).},label={lst:cypher-pass2}]
// Target dikenal sebagai Concept
MATCH (src:Concept {name: $source, grade: $grade})
MATCH (tgt:Concept {name: $target, grade: $grade})
MERGE (src)-[r:{REL_TYPE} {relation_key: $relation_key}]->(tgt)
SET r += $props;

// Target belum dikenal: buat ConceptTarget lalu hubungkan
MERGE (t:ConceptTarget {name: $name, grade: $grade})
SET t.source = 'consensus', t.graph_role = 'target';
MATCH (src:Concept {name: $source, grade: $grade})
MATCH (tgt:ConceptTarget {name: $target, grade: $grade})
MERGE (src)-[r:{REL_TYPE} {relation_key: $relation_key}]->(tgt)
SET r += $props;

// Relasi antar-bab (PRASYARAT / MEMPERSIAPKAN)
MATCH (src:Chapter {name: $source, grade: $grade})
MATCH (tgt:Chapter {name: $target, grade: $grade})
MERGE (src)-[r:{REL_TYPE} {relation_key: $relation_key}]->(tgt)
SET r.description = $description, r.relation_type = $relation_type,
    r.concept_links_json = $concept_links_json,
    r.source = 'consensus', r.graph_role = 'structure';
\end{lstlisting}

\subsubsection{\textit{Writeback} Completion --- Relasi Lintas-Buku}

Setelah tahap \textit{completion}, \textit{edge} lintas-buku berprefiks
\texttt{LINTAS\_BUKU\_} ditulis kembali ke Neo4j dengan
\texttt{graph\_role = 'semantic'}, sehingga dapat dibedakan dari relasi
struktural maupun intra-buku saat analisis pre/post-KGC. Langkah ini bukan
bagian dari konstruksi konsensus dua-\textit{pass}, melainkan \textit{writeback}
hasil KGC.

\begin{lstlisting}[style=cypherbox,caption={Penulisan edge lintas-buku hasil completion.},label={lst:cypher-completion-writeback}]
// Pastikan konsep target ada
MERGE (tgt:Concept {name: $name, grade: $grade})
SET tgt.source = coalesce(tgt.source, 'consensus'),
    tgt.graph_role = coalesce(tgt.graph_role, 'concept');

// Edge lintas-buku
MATCH (src:Concept {name: $source, grade: $source_grade})
MATCH (tgt:Concept {name: $target, grade: $target_grade})
MERGE (src)-[r:{REL_TYPE} {relation_key: $relation_key}]->(tgt)
SET r.description = $description, r.relation_type = $relation_type,
    r.source = 'consensus', r.graph_role = 'semantic';
\end{lstlisting}

% ==================================================================
\subsection{Kueri Perhitungan Metrik}
\label{lamp:cypher-metrik}

Kueri berikut menghitung metrik yang didefinisikan pada Bab~2. Metrik
AD dihitung pada \textbf{subgraf konsep} --- hanya \textit{node} \texttt{:Concept} dan relasi
langsung antar-konsep --- agar tidak terpengaruh \textit{node} struktural.

\subsubsection{Average Degree (AD)}

\begin{lstlisting}[style=cypherbox,caption={AD pada subgraf konsep.},label={lst:cypher-ad}]
MATCH (m:Concept)
WITH m, COUNT { (m)--(:Concept) } AS degree
RETURN avg(degree) AS average_concept_degree
\end{lstlisting}

\subsubsection{Modularity}

Partisi komunitas ditetapkan tetap berdasarkan mata pelajaran (properti
\texttt{grade} pada \texttt{:Concept}), bukan deteksi komunitas otomatis,
sehingga perubahan $Q$ pre/post-KGC murni mencerminkan bertambahnya
\textit{edge} lintas-mata-pelajaran. Kueri mengimplementasikan bentuk ringkas
berarah $Q = \sum_{c}\left[\frac{e_{cc}}{m} - \frac{k_c^{\text{out}}k_c^{\text{in}}}{m^2}\right]$.

\begin{lstlisting}[style=cypherbox,caption={Modularity berarah dengan partisi tetap = mata pelajaran.},label={lst:cypher-modularity}]
// 1. Total edge berarah antar-konsep (m)
MATCH (:Concept)-[r]->(:Concept)
WITH count(r) AS m

// 2. e_cc (edge internal komunitas) dan k_out per komunitas sumber
MATCH (a:Concept)-[r]->(b:Concept)
WITH m, a.grade AS comm,
     count(r) AS k_out,
     sum(CASE WHEN a.grade = b.grade THEN 1 ELSE 0 END) AS e_cc
WITH m, collect({comm: comm, k_out: k_out, e_cc: e_cc}) AS out_stats

// 3. k_in per komunitas target
MATCH (a:Concept)-[r]->(b:Concept)
WITH m, out_stats, b.grade AS comm, count(r) AS k_in
WITH m, out_stats, collect({comm: comm, k_in: k_in}) AS in_stats

// 4. Gabungkan per komunitas dan jumlahkan Q
UNWIND out_stats AS os
WITH m, os.e_cc AS e_cc, os.k_out AS k_out,
     coalesce([x IN in_stats WHERE x.comm = os.comm | x.k_in][0], 0) AS k_in
RETURN sum( (1.0 * e_cc / m) - (1.0 * k_out * k_in / (m * m)) ) AS modularity
\end{lstlisting}

\subsubsection{Uji Ketahanan dengan Mengeluarkan Relasi \texttt{wrong}}

Untuk menghasilkan kolom ``\textit{Wrong} Dikeluarkan'' pada
Tabel~\ref{tab:struktural-robustness}, seluruh pola pencocokan relasi
\texttt{(:Concept)-[r]-(:Concept)} atau
\texttt{(:Concept)-[r]->(:Concept)} pada kueri AD
dan \textit{Modularity} diberi filter berikut. Node \texttt{:Concept} tetap
dipertahankan, sedangkan relasi berlabel \texttt{wrong} dikeluarkan dari
perhitungan.

\begin{lstlisting}[style=cypherbox,caption={Filter untuk uji ketahanan metrik struktural tanpa relasi \texttt{wrong}.},label={lst:cypher-drop-wrong}]
WHERE coalesce(r.consensus, r.status, '') <> 'wrong'
\end{lstlisting}

\subsubsection{Description Completeness (DC)}

\begin{lstlisting}[style=cypherbox,caption={DC atas node :Concept (deskripsi) dan :Chapter (ringkasan).},label={lst:cypher-dc}]
MATCH (n)
WHERE n:Concept OR n:Chapter
WITH n, CASE WHEN n:Concept THEN n.description
             WHEN n:Chapter THEN n.summary END AS desc
WITH count(n) AS total_entities,
     sum(CASE WHEN desc IS NOT NULL AND trim(desc) <> "" THEN 1 ELSE 0 END)
       AS entities_with_description
RETURN total_entities, entities_with_description,
       toFloat(entities_with_description) / total_entities AS description_completeness
\end{lstlisting}

\subsubsection{Empty SubKonsep Rate (ESR)}

\begin{lstlisting}[style=cypherbox,caption={ESR: proporsi :Subtopic tanpa anak :Concept.},label={lst:cypher-esr}]
MATCH (s:Subtopic)
WITH count(s) AS total_subtopics,
     sum(CASE WHEN NOT (s)-[:HAS_CONCEPT]->(:Concept) THEN 1 ELSE 0 END)
       AS empty_subtopics
RETURN total_subtopics, empty_subtopics,
       toFloat(empty_subtopics) / total_subtopics AS empty_subkonsep_rate
\end{lstlisting}

\subsubsection{Typed Relation Ratio (TRR)}

\begin{lstlisting}[style=cypherbox,caption={TRR: proporsi relasi semantik bertipe spesifik (di luar relasi struktural).},label={lst:cypher-trr}]
WITH ["MENDEFINISIKAN","MENYEBABKAN","MEMUNGKINKAN","MENGATUR","BAGIAN_DARI",
      "TERDIRI_DARI","BERGANTUNG_PADA","BERINTERAKSI_DENGAN","BEREAKSI_DENGAN",
      "MENGHASILKAN","MEMPENGARUHI","TERLETAK_DI","DIFORMULASIKAN_SEBAGAI",
      "DIKATALISIS_OLEH","MEMILIKI_SIFAT","LINTAS_BUKU_SAMA_DENGAN",
      "LINTAS_BUKU_APLIKASI_DARI","LINTAS_BUKU_PRASYARAT_UNTUK",
      "LINTAS_BUKU_MEMPERDALAM","LINTAS_BUKU_BERKAITAN_DENGAN"] AS typed_vocab
MATCH ()-[r]->()
WHERE NOT type(r) IN ["HAS_CHAPTER","HAS_SUBTOPIC","HAS_CONCEPT","NEXT_CHAPTER"]
WITH typed_vocab,
     count(r) AS total_semantic_relations,
     sum(CASE WHEN type(r) IN typed_vocab THEN 1 ELSE 0 END) AS typed_relations
RETURN total_semantic_relations, typed_relations,
       toFloat(typed_relations) / total_semantic_relations AS typed_relation_ratio
\end{lstlisting}

\end{document}
